\documentclass[11pt]{article}
\usepackage[margin=1in]{geometry}
\usepackage{amsmath,amssymb,amsthm}
\usepackage{graphicx}
\usepackage[hidelinks]{hyperref}
\setlength{\emergencystretch}{2em}

\newtheorem{theorem}{Theorem}
\newtheorem{lemma}{Lemma}
\newtheorem{corollary}{Corollary}

\newcommand{\M}{\mathcal M}
\newcommand{\K}{\mathcal K}
\newcommand{\B}{\mathcal B}
\newcommand{\N}{\mathbb N}
\newcommand{\T}{\mathbb T}

\title{Sigma-Unital Coronas Fail the Ball-Covering Property}
\author{Run \texttt{fa\_banach\_001}, lane 19}
\date{June 17, 2026}

\begin{document}
\maketitle

\section*{Source Problem}

The source paper is Jinghao Huang, Karimbergen Kudaybergenov, and Rui Liu,
\emph{The ball-covering property of von Neumann algebras and noncommutative
symmetric spaces}, arXiv:2311.14257.  After proving that a commutative unital
\(C^*\)-algebra \(C(K)\) has BCP/UBCP exactly when \(K\) has a countable
\(\pi\)-basis, the authors ask for a noncommutative version:

\begin{center}
\includegraphics[width=0.92\textwidth]{figures/open_problem_crop.png}
\end{center}

\section*{Statement}

Let \(A\) be a nonunital \(C^*\)-algebra.  Write \(\M(A)\) for its multiplier
algebra and
\[
        Q:\M(A)\longrightarrow \M(A)/A
\]
for the quotient map.

\begin{theorem}
\label{thm:sigma-unital-corona}
If \(A\) is nonunital and \(\sigma\)-unital, then the corona
\[
        \M(A)/A
\]
fails the ball-covering property.
\end{theorem}

\begin{corollary}
\label{cor:stable}
Let \(D\) be a nonzero \(\sigma\)-unital \(C^*\)-algebra and let \(H\) be an
infinite-dimensional separable Hilbert space.  Then
\[
        \M(\K(H)\otimes_{\min} D)\big/(\K(H)\otimes_{\min}D)
\]
fails BCP.
\end{corollary}

\begin{proof}
The algebra \(\K(H)\otimes_{\min}D\) is nonunital and \(\sigma\)-unital, so
Theorem~\ref{thm:sigma-unital-corona} applies.
\end{proof}

\section*{Idea of the Proof}

The proof replaces finite coordinate projections by annuli in a sequential
approximate unit.  Choose a positive contractive approximate unit
\((e_n)\) with \(e_{n+1}e_n=e_n\).  For \(T\in\M(A)\), the quotient norm
\(\|Q(T)\|\) is visible on right tails \(T(1-e_n)\), and then on finite
annuli \(T(e_m-e_n)\).  Given countably many proposed centers, recursively
choose disjoint inner annuli where the assigned center nearly realizes its
quotient norm.  Around each inner annulus put a slightly larger annulus and
choose a scalar phase that is Hahn--Banach far from the local piece of the
center.  The strict sum of the phased outer annuli is a multiplier of norm at
most one.  Its corona class has norm one and avoids every proposed ball
\(B(\xi,\|\xi\|)\).

\section*{Preliminaries}

We use the standard equivalent form of BCP: if a normed space \(Y\) has BCP,
then \(S_Y\) is covered by countably many open balls \(B(y_j,\|y_j\|)\).

Since \(A\) is \(\sigma\)-unital, fix a positive contractive approximate unit
\((e_n)\) satisfying
\[
        e_{n+1}e_n=e_n,\qquad n\ge1.
\]
Passing to later terms gives \(e_m e_n=e_n\) whenever \(m\ge n\).  In
particular the \(e_n\)'s commute and \(e_m-e_n\) is a positive contraction
whenever \(m\ge n\).

\begin{lemma}[quotient norm on approximate-unit tails]
\label{lem:tail}
For every \(T\in\M(A)\),
\[
        \|Q(T)\|=\lim_{n\to\infty}\|T(1-e_n)\|.
\]
\end{lemma}

\begin{proof}
Since \(Te_n\in A\), we have \(\|Q(T)\|\le \|T(1-e_n)\|\) for every \(n\).
Conversely, if \(C\in A\), then \(C(1-e_n)\to0\) in norm, and hence
\[
        \limsup_n \|T(1-e_n)\|
        \le
        \limsup_n \|(T-C)(1-e_n)\|+\limsup_n\|C(1-e_n)\|
        \le \|T-C\|.
\]
Taking the infimum over \(C\in A\) gives the reverse inequality.
\end{proof}

\begin{lemma}[finite annulus recovery]
\label{lem:annulus-recovery}
Let \(T\in\M(A)\), put \(r=\|Q(T)\|\), and let \(N\in\N\) and
\(\varepsilon>0\).  Then there are integers
\[
        N<n<m
\]
such that
\[
        \|T(e_m-e_n)\|>r-\varepsilon.
\]
\end{lemma}

\begin{proof}
By Lemma~\ref{lem:tail}, choose \(n>N\) such that
\[
        \|T(1-e_n)\|>r-\varepsilon/2.
\]
For every \(B\in\M(A)\),
\[
        \|B\|=\sup_m \|B e_m\|.
\]
Indeed, the inequality ``\(\ge\)'' is immediate, while for \(a\in A\) with
\(\|a\|\le1\),
\[
        \|Ba\|=\lim_m\|B e_m a\|\le \sup_m\|B e_m\|,
\]
and \(A\) is an essential ideal in \(\M(A)\).  Applying this to
\(B=T(1-e_n)\), choose \(m>n\) so that
\[
        \|T(1-e_n)e_m\|>r-\varepsilon.
\]
Since \(e_ne_m=e_n\), this is exactly
\[
        \|T(e_m-e_n)\|>r-\varepsilon.
\]
\end{proof}

\begin{lemma}[orthogonal strict sums]
\label{lem:strict-sum}
Let
\[
        a_k<n_k<m_k<b_k<a_{k+1},\qquad k\ge1,
\]
and define
\[
        g_k=e_{m_k}-e_{n_k},\qquad h_k=e_{b_k}-e_{a_k}.
\]
Then the \(h_k\)'s are pairwise orthogonal positive contractions,
\[
        h_k g_k=g_k,\qquad h_jg_k=0\quad(j\ne k),
\]
and for every bounded scalar sequence \((\lambda_k)\), the series
\[
        z=\sum_{k=1}^{\infty}\lambda_k h_k
\]
converges strictly to an element of \(\M(A)\) with
\[
        \|z\|\le \sup_k|\lambda_k|,
        \qquad
        z g_k=\lambda_k g_k.
\]
\end{lemma}

\begin{proof}
All identities follow from \(e_me_n=e_{\min(m,n)}\) for \(m,n\) in the chosen
subsequence.  For example, if \(a_k<n_k<m_k<b_k\), then
\[
        (e_{b_k}-e_{a_k})(e_{m_k}-e_{n_k})=e_{m_k}-e_{n_k}=g_k.
\]
If the intervals are disjoint, the same computation gives \(h_jh_k=0\) and
\(h_jg_k=0\) for \(j\ne k\).

For strict convergence, let \(s_N=\sum_{k\le N}\lambda_kh_k\).  The partial
sums are uniformly bounded by \(\sup_k|\lambda_k|\), because the \(h_k\)'s are
orthogonal positive contractions.  If \(a\in A\), then the tail
\(\sum_{k>N}\lambda_kh_k\) is annihilated by \(e_{a_{N+1}}\), so
\[
        \Big\|\Big(\sum_{k>N}\lambda_kh_k\Big)a\Big\|
        =
        \Big\|\Big(\sum_{k>N}\lambda_kh_k\Big)(a-e_{a_{N+1}}a)\Big\|
        \le
        \sup_k|\lambda_k|\,\|a-e_{a_{N+1}}a\|\to0.
\]
The same argument on the other side gives
\[
        \Big\|a\Big(\sum_{k>N}\lambda_kh_k\Big)\Big\|\to0.
\]
Thus the series is strictly convergent in \(\M(A)\), with the stated norm
bound and block action.
\end{proof}

\section*{Proof of the Theorem}

\begin{proof}[Proof of Theorem~\ref{thm:sigma-unital-corona}]
Let
\[
        \xi_i=Q(T_i),\qquad i\ge1,
\]
be an arbitrary countable family of centers in \(\M(A)/A\), with
\(T_i\in\M(A)\).  Put \(r_i=\|\xi_i\|\).  Add one auxiliary multiplier
\[
        T_0=1_{\M(A)},\qquad r_0=1.
\]
Choose a map \(\pi:\N\to\{0,1,2,\ldots\}\) such that every value occurs
infinitely often.

We recursively choose integers
\[
        a_k<n_k<m_k<b_k<a_{k+1}
\]
and phases \(\lambda_k\in\T\).  Suppose the previous choices have been made,
and put \(i=\pi(k)\).  Choose \(a_k>b_{k-1}\), with the convention
\(b_0=0\).  By Lemma~\ref{lem:annulus-recovery}, choose
\[
        a_k<n_k<m_k
\]
so that, with \(g_k=e_{m_k}-e_{n_k}\),
\[
        \|T_i g_k\|>r_i-\frac1k.
\]
Choose \(b_k>m_k\), and later choose \(a_{k+1}>b_k\).

Now choose \(\lambda_k\in\T\) so that
\[
        \|\lambda_k g_k-T_i g_k\|\ge \|T_i g_k\|.
\]
This is just the Hahn--Banach sign trick: if \(T_i g_k\ne0\), take a norm-one
functional \(\varphi\) with \(\varphi(T_i g_k)=\|T_i g_k\|\), and choose
\(\lambda_k\in\T\) with
\(\operatorname{Re}(\lambda_k\varphi(g_k))\le0\).  If \(T_i g_k=0\), any
\(\lambda_k\) works.

Let
\[
        h_k=e_{b_k}-e_{a_k}.
\]
By Lemma~\ref{lem:strict-sum}, the series
\[
        z=\sum_{k=1}^{\infty}\lambda_k h_k
\]
converges strictly to a multiplier with \(\|z\|\le1\), and
\[
        z g_k=\lambda_k g_k.
\]

We first show that \(\|Q(z)\|=1\).  The upper bound follows from \(\|z\|\le1\).
Let \(C\in A\).  Take \(k\to\infty\) along the infinite set
\(\{k:\pi(k)=0\}\).  For these \(k\)'s,
\[
        \|g_k\|=\|T_0g_k\|>1-\frac1k.
\]
Moreover \(n_k\to\infty\), so
\[
        \|C g_k\|
        =
        \|C(1-e_{n_k})e_{m_k}\|
        \le
        \|C(1-e_{n_k})\|\longrightarrow0.
\]
Therefore
\[
        \|z-C\|
        \ge
        \|(z-C)g_k\|
        \ge
        \|z g_k\|-\|C g_k\|
        =
        \|g_k\|-\|C g_k\|,
\]
and letting \(k\to\infty\) along \(\pi(k)=0\) gives \(\|z-C\|\ge1\).  Taking
the infimum over \(C\in A\), we get \(\|Q(z)\|\ge1\).

Now fix \(i\ge1\) and \(C\in A\).  Let \(k\to\infty\) along the infinite set
\(\{k:\pi(k)=i\}\).  Since \(n_k\to\infty\), \(\|Cg_k\|\to0\).  Hence
\[
        \|z-T_i-C\|
        \ge
        \|(z-T_i-C)g_k\|
        \ge
        \|\lambda_k g_k-T_ig_k\|-\|Cg_k\|
        \ge
        \|T_ig_k\|-\|Cg_k\|
        >
        r_i-\frac1k-\|Cg_k\|.
\]
Letting \(k\to\infty\) along this subsequence gives
\[
        \|z-T_i-C\|\ge r_i.
\]
Taking the infimum over \(C\in A\), we obtain
\[
        \|Q(z)-\xi_i\|\ge r_i=\|\xi_i\|.
\]

Thus \(Q(z)\) is a unit vector in \(\M(A)/A\) outside every open ball
\[
        B(\xi_i,\|\xi_i\|).
\]
Since the countable family of centers was arbitrary, \(\M(A)/A\) fails BCP.
\end{proof}

\section*{Scope and Relation to Earlier Packets}

The previous stabilized-corona packet proved the result for
\[
        (\B(H)\otimes D)/(\K(H)\otimes D)
\]
with unital coefficient algebra \(D\), using finite-rank coordinate
projections and partial isometries.  The present theorem is both simpler and
broader for multiplier coronas: it applies to every nonunital
\(\sigma\)-unital \(C^*\)-algebra \(A\).  In particular it covers
\[
        A=\K(H)\otimes_{\min}D
\]
for every nonzero \(\sigma\)-unital coefficient algebra \(D\), unital or not.

The theorem is still a partial answer to the source question.  It does not
characterize arbitrary \(C^*\)-algebras with BCP, and it does not address
coronas of non-\(\sigma\)-unital ideals.

\section*{Verification and Search Bounds}

The proof is self-contained modulo the standard fact that a \(\sigma\)-unital
\(C^*\)-algebra has a positive contractive approximate unit
\((e_n)\) satisfying \(e_{n+1}e_n=e_n\).  The run indexes were searched for
arXiv:2311.14257 together with BCP, UBCP, corona, multiplier algebra,
\(\sigma\)-unital, approximate unit, Calkin, and stable multiplier keywords.
Existing run packets covered product quotients, reduced products, continuous
fields, unital stabilized coronas, and a Hilbert-module block criterion, but
no packet contained this approximate-unit annulus theorem.

External literature checks during this push found arXiv:2412.07137, which
proves that several Calkin algebras fail BCP, and arXiv:2111.04921, which
gives commutative/vector-valued BCP stability results and warns that BCP is
not hereditary even for 1-complemented subspaces.  Those papers are consistent
with the present theorem but do not state the sigma-unital multiplier-corona
annulus result used here.

\section*{Human Review Focus}

Review Lemma~\ref{lem:annulus-recovery} and Lemma~\ref{lem:strict-sum}.
These are the two structural points: quotient norms must be recoverable on
finite approximate-unit annuli, and the phased outer annuli must define a
strict multiplier acting as the chosen phase on the inner annulus.  Once those
are checked, the final diagonal escape is formal.

\begin{thebibliography}{9}

\bibitem{HKL}
Jinghao Huang, Karimbergen Kudaybergenov, and Rui Liu,
\emph{The ball-covering property of von Neumann algebras and noncommutative
symmetric spaces},
arXiv:2311.14257, 2023.

\bibitem{LLLZ}
Minzeng Liu, Rui Liu, Jimeng Lu, and Bentuo Zheng,
\emph{Ball covering property from commutative function spaces to
non-commutative spaces of operators},
arXiv:2111.04921, 2021.

\bibitem{BLS}
Qiyao Bao, Rui Liu, and Jie Shen,
\emph{The ball-covering property of non-commutative spaces of operators on
Banach spaces},
arXiv:2412.07137, 2024.

\end{thebibliography}

\end{document}
